% ============================================================================
%  Part VI: Security Analysis
% ============================================================================
\section{Security Analysis}
\label{sec:security}

This part lays out the threat model PQID is designed against, the cryptographic primitives that back it, and the attacks it does \emph{not} defend against.

\subsection{Threat model}

PQID is designed to provide three security properties:

\begin{enumerate}
  \item \textbf{Uniqueness.} The probability that two independently-generated PQIDs collide is vanishingly small under both classical and quantum adversaries.
  \item \textbf{Unguessability.} An attacker who has seen many PQIDs cannot guess a PQID that has not yet been generated, or a PQID generated by a different system, with non-negligible probability.
  \item \textbf{Tamper-evidence under outer authentication.} When a PQID is wrapped in a signed envelope (e.g., a JWT \citep{rfc7519} or PASETO \citep{paseto} token), modifying the PQID invalidates the signature. PQID itself does not provide authentication; that is the envelope's job.
\end{enumerate}

PQID is \emph{not} designed to provide:

\begin{itemize}
  \item \textbf{Authentication of the producer.} Anyone with a CSPRNG can forge a valid-looking PQID. Use a signed envelope (PASETO, JWT, custom MAC) if you need to verify the producer.
  \item \textbf{Integrity of the identifier itself.} PQID has no checksum, MAC, or hash. A bit-flipped PQID may decode to another valid PQID. If integrity matters, the PQID belongs inside a signed or HMAC'd \citep{rfc2104} container that protects it.
  \item \textbf{Confidentiality of the timestamp.} The 48-bit timestamp is plaintext. An attacker who sees a PQID learns approximately when it was generated. If timestamp confidentiality matters, encrypt the entire identifier.
\end{itemize}

\subsection{Quantum-relevant attacks}

\subsubsection{Grover's algorithm \citep{grover1996}}

Grover's algorithm searches an unstructured space of $N$ elements in $O(\sqrt{N})$ operations. For an $n$-bit random field, this means finding a single specific value (a targeted second preimage) costs $2^{n/2}$ quantum operations instead of $2^n$ classical operations. For PQID's 256-bit random field, that's $2^{128}$ --- astronomically out of reach for any plausible quantum computer.

\subsubsection{Brassard--H{\o}yer--Tapp \citep{bht1998}}

BHT extends Grover to find collisions: any two values in the space that share an output. For an $n$-bit space, BHT finds a collision in $O(2^{n/3})$ operations. For PQID's 256-bit field, that's $2^{85}$ --- large but not as comfortable as the 128-bit second-preimage bound. The relevant point is that for \emph{identifier} use cases, second-preimage is the threat model that matters: an attacker wants to collide with a specific known PQID, not find any two PQIDs that collide. For audit-log integrity, second-preimage is the right bound, and 128 bits is the target.

\subsection{Birthday-bound math at scale}

The classical birthday bound for an $n$-bit random field is approximately $\sqrt{\pi/2} \cdot 2^{n/2}$ generations before an expected collision. For PQID's 256-bit tail:

\begin{equation}
\sqrt{\pi/2} \cdot 2^{128} \approx 1.253 \cdot 3.40 \times 10^{38} \approx 4.27 \times 10^{38} \text{ generations}
\end{equation}

At a sustained rate of $10^9$ IDs per second --- a conservative estimate for the largest systems imaginable in the late 2020s --- this is approximately $1.35 \times 10^{22}$ years. The current age of the universe is approximately $1.4 \times 10^{10}$ years, so PQID's classical birthday bound is roughly $10^{12}$ times the universe's current age, or one trillion times.

For comparison:

\begin{center}
\renewcommand{\arraystretch}{1.15}
\begin{tabular}{lrl}
\toprule
\textbf{Format} & \textbf{Random bits} & \textbf{Years to first classical collision @ $10^9$/s} \\
\midrule
PQID    & 256 & $1.35 \times 10^{22}$ \\
KSUID   & 128 & $1.4 \times 10^{10}$  \\
ULID    &  80 & $1.0 \times 10^{2}$  \\
UUIDv7  &  74 & $5.0 \times 10^{0}$  \\
UUIDv4  & 122 & $4.4 \times 10^{8}$  \\
\bottomrule
\end{tabular}
\end{center}

UUIDv7's 5-year classical birthday horizon at $10^9$/s is the number that motivates the AI-scale discussion in Part \ref{sec:ai-scale}: at million-node training-cluster scales, even classical collision math becomes a real concern for UUIDv7, before any quantum considerations enter.

\subsection{Side channels}

PQID generators leak information about generation rate via the counter. An attacker who can observe many PQIDs from a single generator can estimate the generation rate from the rate at which the counter advances within a millisecond. This is generally not a meaningful information leak (the timestamp already gives away when the ID was generated, and the counter advances by at most 16 per ID), but it is worth noting.

The HLC variant additionally leaks information about peer observations through its logical clock advance. An HLC PQID whose logical time is far ahead of the local wall clock indicates the generator recently observed a PQID from a peer with a faster-moving clock.

If side-channel resistance matters, encrypt the PQID before publication.

\subsection{Random source quality}

PQID's security guarantees rest entirely on the quality of the underlying CSPRNG. All reference implementations use the operating system's CSPRNG (\code{getrandom(2)} on Linux, \code{arc4random\_buf(3)} on macOS/BSD, \code{BCryptGenRandom} on Windows). These are designed for cryptographic use and are continuously vetted by their respective security teams.

Implementations MUST NOT use language-level non-cryptographic random sources (\code{rand()}, \code{Math.random()}, \code{random.random()}) for the PQID random tail. Such sources are predictable and would catastrophically reduce PQID's effective security.

